\chapter{Commodities}
\label{ch:commodities}
\chaptermark{Commodities}

Commodities are physical goods with storage costs, seasonality, and
the tightest link to the real economy.  A barrel of crude has to sit
somewhere and somebody pays the rent, which warps every derivative
price.  This chapter covers futures-curve mechanics, the
cost-of-carry formula extended with \emph{storage cost} $u$ and
\emph{convenience yield} $y$ (the hidden dividend earned by whoever
has physical inventory when a refinery runs dry), and Schwartz's
one-factor mean-reverting model.

\begin{backgroundbox}[Prerequisites]
From earlier chapters:
\begin{itemize}
  \item \Cref{ch:arbitrage_zoo}: the cash-and-carry no-arbitrage
    argument for forward prices on storable assets, $F = S e^{\rrf T}$
    for a pure-investment asset.
  \item \Cref{ch:options_payoffs}: call and put payoffs, risk-neutral
    expectation, the Black--Scholes pricing machinery.
  \item \Cref{ch:fixed_income_foundations}: the discount factor
    $D(t) = e^{-\rrf t}$ under a flat zero curve, used to PV future
    cash flows.
\end{itemize}
If the cost-of-carry argument of \cref{ch:arbitrage_zoo} feels
hazy, re-read that chapter first --- this one builds directly on it.
\end{backgroundbox}


%% ================================================================
\section{Futures curves: contango and backwardation}
\label{sec:contango_backwardation}
%% ================================================================

A \emph{futures curve} (or \emph{forward curve}) is the function
$T \mapsto F(T)$ mapping each maturity to today's price for delivery
at $T$.  For a stock it is an upward-sloping exponential dictated by
the risk-free rate.  For a commodity it has shape: it can slope up,
slope down, bend, or wiggle seasonally, encoding everything the
market believes about storage, supply, demand, and inventory.

\begin{definition}[Contango]
\index{contango|textbf}
A futures curve is in \emph{contango} if longer-dated futures trade
above shorter-dated futures:
\[
  T_2 > T_1 \implies F(T_2) > F(T_1).
\]
The curve slopes upward.
\end{definition}

\begin{definition}[Backwardation]
\index{backwardation|textbf}
A futures curve is in \emph{backwardation} if longer-dated futures
trade below shorter-dated futures:
\[
  T_2 > T_1 \implies F(T_2) < F(T_1).
\]
The curve slopes downward.
\end{definition}

Real curves can be contangoed at the front and backwardated at the
back, or vice versa.  Traders therefore speak of the \emph{front-
month spread} $F(T_2)-F(T_1)$ between the two nearest contracts as
the local measure of curve shape: positive means local contango,
negative means local backwardation.

\subsection{A worked contango curve}

Consider West Texas Intermediate (WTI) crude with spot $\Spx = \$70$
per barrel and a net cost-of-carry $r + u - y$ varying mildly with
maturity.  A plausible contangoed curve:

\begin{center}
\begin{tabular}{lcc}
\toprule
Maturity $T$ & Net carry $r+u-y$ & Futures price $F(T)$ \\
\midrule
1 month  & $3.0\%$ & \$70.18 \\
3 months & $3.5\%$ & \$70.62 \\
6 months & $4.0\%$ & \$71.41 \\
12 months & $5.0\%$ & \$73.59 \\
\bottomrule
\end{tabular}
\end{center}

Each row is $F(T) = 70 \cdot e^{(r+u-y)T}$.  The 12-month future
trades \$3.59 above spot.  This is the \emph{normal} shape for
storable commodities with ample inventories and no scarcity
premium.

\subsection{A worked backwardation curve}

Now suppose a refinery outage spikes demand for \emph{prompt}
barrels.  Spot rises but distant futures do not.  Model this as
large negative net carry ($y$ dominates $r + u$):

\begin{center}
\begin{tabular}{lcc}
\toprule
Maturity $T$ & Net carry $r+u-y$ & Futures price $F(T)$ \\
\midrule
1 month  & $-8.0\%$ & \$69.54 \\
3 months & $-6.0\%$ & \$68.96 \\
6 months & $-4.0\%$ & \$68.61 \\
12 months & $-2.0\%$ & \$68.61 \\
\bottomrule
\end{tabular}
\end{center}

The curve slopes downward at the front: ``hold a barrel today and
it will be worth more than a promise of one in a month.''

\begin{keyfactbox}[Contango vs backwardation]
\begin{itemize}
  \item \textbf{Contango}: $F(T)$ increases in $T$.  Normal for
    storable commodities with loose inventories and no scarcity.
    \emph{Long futures bleeds}: as time passes the contract rolls
    down the curve toward the lower spot price.
  \item \textbf{Backwardation}: $F(T)$ decreases in $T$.  Typical
    for supply-constrained or expiring stockpile regimes.
    \emph{Long futures earns a roll yield}: the contract rolls up
    the curve toward the higher spot price as it ages.
\end{itemize}
The sign of the curve slope is the single most-watched indicator in
a commodity dealing room.
\end{keyfactbox}

\begin{intuitionbox}
Backwardated = market pays you to hold physical now: you earn an
implicit dividend (the \emph{convenience yield}) that shows up
nowhere on P\&L but reveals itself in the curve's shape.  Contango =
market pays you to defer; warehouses are full and holding costs
storage and financing with no offsetting benefit.
\end{intuitionbox}

\subsection{Roll yield and commodity index returns}
\label{subsec:roll_yield}

A commodity index fund (e.g.\ an ETF tracking the Bloomberg
Commodity Index or S\&P GSCI) does not own physical barrels.  It
holds front-month futures and on a scheduled monthly ``roll window''
sells the expiring contract and buys the next.  The fund's return
decomposes into three pieces:
\[
  r_{\text{index}} \;=\; \underbrace{r_{\text{spot}}}_{\text{spot return}}
    \;+\; \underbrace{r_{\text{roll}}}_{\text{roll yield}}
    \;+\; \underbrace{r_{\text{collateral}}}_{\text{T-bill rate}}.
\]
\emph{Collateral yield} comes from investing the notional cash in
short T-bills.  \emph{Spot return} is the percentage change in the
commodity spot (proxied by the nearest-futures price).  \emph{Roll
yield} is the kicker:
\[
  r_{\text{roll}} \approx \frac{F(T_1) - F(T_2)}{F(T_2)},
\]
negative in contango (sell low, buy high), positive in
backwardation.  Contango is a persistent drag on index returns: the
crude-oil sub-index lost 5--10\% per year to negative roll yield
during the 2014--2016 contango.

\begin{intuitionbox}
An index fund tracking a contangoed commodity pays the carry
embedded in the term structure every month.  This is the biggest
reason naive ``long commodity'' strategies have underperformed over
the last two decades: not because spot went nowhere, but because
they paid carry continuously.
\end{intuitionbox}

\begin{pitfallbox}[Contango $\ne$ expected price path]
A common beginner error reads a contangoed curve as ``the market
expects prices to rise.''  It does not.  $F(T)$ is a risk-adjusted
delivery price set by carry economics, not a forecast of $\Spx_T$.
The expected spot $\E[S_T]$ can equal, exceed, or fall below $F(T)$
depending on the risk premium.  Keynes's \emph{normal backwardation}
theory argues producers hedging forward supply push
$F(T) < \E[S_T]$, making backwardation the structurally ``normal''
state.  Do not confuse curve shape with directional view.
\end{pitfallbox}


%% ================================================================
\section{Cost-of-carry with storage and convenience yield}
\label{sec:cost_of_carry_commodities}
%% ================================================================

\Cref{ch:arbitrage_zoo} derived $F = \Spx e^{\rrf T}$ for a pure-
investment asset by cash-and-carry replication.  For commodities
this breaks in two ways.  First, holding the physical barrel is not
free: you pay \emph{storage costs} (tank rental, insurance, boil-
off for LNG).  Second, holding the barrel confers a \emph{benefit}
--- if your refinery runs dry you can feed it, whereas a futures
contract cannot be cracked into gasoline.  That benefit is the
\emph{convenience yield}.

\begin{definition}[Storage cost rate]
\index{storage cost|textbf}
The \emph{storage cost rate} $u$ is the continuously-compounded
rate at which a holder of one unit of physical commodity pays to
store it.  A typical crude storage cost of \$0.50/barrel/month on a
\$70 barrel corresponds to $u \approx 8.6\%$ per year; realised
values are smaller because tanks are not full.  The modelling
idealisation is continuous payment proportional to spot value.
\end{definition}

\begin{definition}[Convenience yield]
\index{convenience yield|textbf}
The \emph{convenience yield} $y$ is the continuously-compounded
implicit benefit rate of holding physical inventory versus a
futures contract on it.  It captures the optionality to meet a
supply shock, avoid a stockout, or feed a downstream process, and
is applied per-year to the spot value.
\end{definition}

\subsection{Deriving the cost-of-carry formula}

Extend the \cref{ch:arbitrage_zoo} cash-and-carry argument.  At time
0 choose between two strategies delivering one unit of commodity at
$T$:

\medskip
\noindent\textbf{Strategy A: hold futures.} Buy one long futures
contract at price $F$.  Invest $F e^{-\rrf T}$ in the risk-free
asset; at $T$ the bond pays $F$, you pay $F$ to take delivery.
Net: one unit of commodity at $T$, cost today $F e^{-\rrf T}$.

\medskip
\noindent\textbf{Strategy B: cash-and-carry.} Borrow $\Spx$ at rate
$\rrf$, buy spot, pay storage continuously at rate $u$ applied to
the value of the carried physical, and enjoy convenience yield $y$
also applied to the value of the carried physical.  At $T$:
\begin{itemize}
  \item cost of carried physical (spot + storage - convenience)
    grows at rate $\rrf + u - y$, so the equivalent time-$T$
    accumulated cost is $\Spx e^{(\rrf+u-y)T}$;
  \item the physical unit itself is delivered to you.
\end{itemize}
Net: one unit of commodity at $T$, accumulated cost
$\Spx e^{(\rrf+u-y)T}$.

\medskip
Both strategies replicate ``one unit of commodity at $T$.''
No-arbitrage equates their costs: $F = \Spx e^{(\rrf+u-y)T}$.

\begin{keyfactbox}[Commodity forward pricing]
For a storable commodity with continuously-compounded storage cost
rate $u$ and convenience yield $y$,
\[
  \boxed{\; F(T) = \Spx\, e^{(\rrf + u - y)\,T}. \;}
\]
\begin{itemize}
  \item $u > y$ (loose inventories, storage dominates):
    $F > \Spx e^{\rrf T}$, curve in contango above the pure
    financing curve.
  \item $y > u + \rrf$ (tight inventories, convenience dominates):
    $F < \Spx$, curve in backwardation below spot.
  \item Equities and currencies are special cases: set $u=0$ and
    $y$ equal to the dividend or foreign rate.
\end{itemize}
The sign of $\rrf + u - y$ is the sign of the curve slope at $T$
--- exactly the continuously-compounded front-to-back roll rate.
\end{keyfactbox}

\subsection{Worked example: crude oil forward}

Let spot crude $\Spx = \$80$, risk-free rate $\rrf = 5\%$, storage
cost $u = 2\%$, convenience yield $y = 4\%$, horizon $T = 1$ year.
Net carry:
\[
  \rrf + u - y = 0.05 + 0.02 - 0.04 = 0.03.
\]
Forward price:
\[
  F(1) = 80 \cdot e^{0.03 \cdot 1} = 80 \cdot 1.03045 = \$82.44.
\]
(Verified: $80 e^{0.03} = 82.4364$.) Mild contango: the forward
trades \$2.44 above spot as storage plus financing slightly exceeds
convenience.  If $y$ rose to $8\%$, net carry flips to $-1\%$ and
the one-year forward drops to $80 e^{-0.01} = \$79.20$
(backwardation).

\subsection{Where does convenience yield come from? The theory of storage}
\label{subsec:theory_of_storage}

The convenience yield has economic content beyond a fudge factor.
The Kaldor--Working \emph{theory of storage} provides the clearest
picture.  A firm chooses each period whether to hold one extra
barrel.  The benefit is avoiding a stockout: if demand spikes or
supply is disrupted, the firm feeds its refinery from inventory
rather than scrambling in the spot market at punitive prices.

Formally, let $C(I)$ be the expected cost of a stockout per unit
time when inventory level is $I$ (decreasing and convex in $I$:
more inventory, lower stockout cost).  The marginal convenience
benefit of the $I$-th barrel is $-C'(I)$, which is positive and
decreasing in $I$.  In competitive equilibrium, a firm holds
inventory up to the point where this marginal benefit equals the
marginal cost of storage:
\[
  \underbrace{-C'(I^*)}_{\text{marginal convenience}}
    \;=\; \underbrace{\rrf + u}_{\text{cost of carry}}
    \;-\; \underbrace{\E[\Delta \ln S]}_{\text{expected spot drift}}.
\]
The left-hand side is the convenience yield $y$.  Low inventories
(near stockout) make $-C'(I)$ large, $y$ is large, and the curve
backwardates.  Full tanks make $-C'(I) \approx 0$, so $y \approx 0$
and the curve defaults to contango driven by $\rrf + u$.

This explains the empirical regularity: \emph{curve shape is a
near-monotonic function of the inventory-to-use ratio}.  The
``storage arbitrage line'' plots $y$ against days-of-supply as a
smooth downward-sloping curve that shifts seasonally.

\begin{intuitionbox}
Convenience yield is the \emph{shadow price of an empty tank}:
large when inventories are scarce (the next barrel has optionality
to be sold into a shortage at a huge mark-up), near zero when tanks
are full.  A curve move from contango to backwardation is almost
always a statement about \emph{inventory tightening}, not about
expected future price paths.
\end{intuitionbox}

\subsection{Inferring convenience yield from the curve}

$y$ is not quoted; the market quotes $\{F(T_i)\}$ and
$\{\rrf(T_i)\}$, and storage $u$ is known from physical contracts.
Inverting cost-of-carry,
\[
  y(T) = \rrf(T) + u(T) - \frac{1}{T} \ln\frac{F(T)}{\Spx}.
\]
This \emph{implied convenience yield} is treated as a market-
observable state variable --- the commodity analogue of extracting
dividend yields from equity forward markets.

\subsection{Calendar spreads and the curve as a tradeable instrument}
\label{subsec:calendar_spreads}

The futures curve is itself a tradeable object.  A \emph{calendar
spread} is a simultaneous long-one-maturity / short-another
position:
\[
  \text{CL}\,\text{Dec}\text{--}\text{Jun spread}
    \;=\; F(T_{\text{Dec}}) - F(T_{\text{Jun}}).
\]
A long spread profits if the curve steepens; short profits if it
flattens.  Spread volatilities are an order of magnitude below
outright futures vols because both legs share the systematic level
risk --- what remains is \emph{curve-shape} risk, driven by
inventories and convenience yield.

Exchanges list calendar spreads as explicit products (e.g.\ ICE
Brent ``BZU'' month-to-month spreads) with their own order books,
because synthesising them from outright legs carries execution risk.
Inter-commodity spreads --- WTI vs Brent (``trans-Atlantic''),
heating oil vs crude (``heat crack''), gasoline-crude (``crack
3:2:1'') --- dominate physical refinery hedging.  The
\emph{crack spread} is the refiner's margin: the price of
refined products (gasoline, heating oil) minus the crude
input, e.g.\ the 3:2:1 crack $=\tfrac13(2\,\text{RBOB} +
\text{HO}) - \text{WTI}$.  The \emph{spark spread} is the
analogous power-generator margin: the price of electricity
minus the fuel cost of generating it, typically
$\text{Power} - H_R\cdot \text{NatGas}$ where $H_R$ is the
plant's heat rate (MMBtu per MWh).

\begin{keyfactbox}[Relation to convenience yield]
The one-period log calendar spread equals the forward-forward net
cost-of-carry:
\[
  \ln F(T_2) - \ln F(T_1)
    = (\rrf + u - y)(T_2 - T_1) \quad\text{(flat-carry case)}.
\]
Calendar-spread trades are \emph{directional bets on changes in
convenience yield $y$}.  When inventory draws faster than expected,
$y$ rises at the front, $F(T_1)$ rises relative to $F(T_2)$, and the
spread $F(T_2) - F(T_1)$ compresses.
\end{keyfactbox}

\begin{gsbox}[title={Curve construction on the commodities desk}]
A Goldman commodities Strat maintains three curves per product:
(i) the \emph{futures curve} $F(T)$ bootstrapped from listed
contracts (e.g.\ 36 monthly CL contracts on NYMEX); (ii) the
\emph{spot/prompt curve} for deliverable physical plus near-term
basis; (iii) the \emph{implied convenience yield curve} $y(T)$
backed out from the above and the Treasury zero curve.  Equities
need only spot and the dividend curve; here $y(T)$ is where all the
commodity-specific alpha and risk live --- it moves sharply with
inventory data, OPEC announcements, and pipeline disruptions.
Option Strats shock $y$ by $\pm 100$ bp and mark the P\&L; it is a
first-class risk factor alongside the futures price.
\end{gsbox}


%% ================================================================
\section{The Schwartz one-factor model}
\label{sec:schwartz_model}
%% ================================================================

Cost-of-carry is a relation between spot and forward; it is silent
on spot dynamics.  For risk management, options pricing, and curve
evolution we need a model of $\Spx_t$.  Geometric Brownian motion
fits commodities poorly: it implies unbounded drift, whereas crude,
copper, and natural gas all \emph{mean-revert} to a long-run
production-cost level.  When WTI hit $-\$37$ in April 2020 or \$140
in 2008 it did not stay there; supply and demand pulled it back.

\citet{schwartz1997stochastic} proposed the canonical one-factor
mean-reverting model: log-spot follows an Ornstein--Uhlenbeck
process, so spot follows a geometric mean-reverting SDE.

\begin{definition}[Schwartz one-factor model]
\index{Schwartz model|textbf}
Let $X_t \defeq \ln \Spx_t$.  Under the real-world measure,
\begin{equation}
  dX_t = \kappa\,(\alpha - X_t)\,dt + \volat\,dW_t,
  \label{eq:schwartz_ou_log}
\end{equation}
where $\kappa > 0$ is the mean-reversion speed, $\alpha$ the
long-run log-mean, and $\volat > 0$ the volatility of log-spot.
Equivalently, in the spot variable,
\begin{equation}
  d\Spx_t = \kappa\,\bigl(\mu - \ln \Spx_t\bigr)\,\Spx_t\,dt
          + \volat\,\Spx_t\,dW_t,
  \label{eq:schwartz_spot}
\end{equation}
with $\mu = \alpha + \volat^2/(2\kappa)$ after the Itô correction
from log to level.
\end{definition}

$X_t$ is Gaussian Ornstein--Uhlenbeck: reverts to $\alpha$ at rate
$\kappa$ with steady-state variance $\volat^2/(2\kappa)$ and shock
half-life $\ln 2 / \kappa$.  Typical WTI calibration: $\kappa
\approx 1$/yr (half-life $\approx 8$ months), $\volat \approx 30\%$.
Electricity reverts in days, not months.

\subsection{The futures price under Schwartz}

Under the risk-neutral measure $\Qm$, log-spot still follows an OU
process but with risk-adjusted long-run mean $\alpha^*$:
\[
  dX_t = \kappa\,(\alpha^* - X_t)\,dt + \volat\,d\widetilde W_t.
\]
The futures price is $F(t,T) = \E^\Qm[\Spx_T \mid \mathcal F_t]$.
$X_T$ is Gaussian under $\Qm$ with
\[
  \E^\Qm[X_T \mid \mathcal F_t] = e^{-\kappa(T-t)} X_t
    + (1 - e^{-\kappa(T-t)})\alpha^*,
\]
\[
  \Var^\Qm[X_T \mid \mathcal F_t]
    = \frac{\volat^2}{2\kappa}\bigl(1 - e^{-2\kappa(T-t)}\bigr),
\]
and $\Spx_T = e^{X_T}$ is lognormal, we use
$\E[e^{X_T}] = \exp(\E[X_T] + \tfrac12 \Var[X_T])$ to obtain the
Schwartz closed-form futures price.

\begin{keyfactbox}[Schwartz closed-form futures]
Under the one-factor Schwartz model with risk-neutral long-run
log-mean $\alpha^*$, reversion speed $\kappa$, and volatility
$\volat$, the time-$t$ futures price for delivery at $T$ is
\begin{equation}
  F(t,T) = \exp\!\Bigl[
    e^{-\kappa(T-t)} \ln \Spx_t
    + (1 - e^{-\kappa(T-t)})\,\alpha^*
    + \frac{\volat^2}{4\kappa}\bigl(1 - e^{-2\kappa(T-t)}\bigr)
  \Bigr].
  \label{eq:schwartz_F}
\end{equation}
As $T \to \infty$, $F(t,T) \to \exp(\alpha^* + \volat^2/(4\kappa))$,
a constant independent of current spot --- the model's long-run
equilibrium price.
\end{keyfactbox}

\begin{intuitionbox}
The long end of the futures curve is anchored at the equilibrium
level $\alpha^*$, independent of spot.  Only the front moves with
spot --- the empirical stylised fact: when WTI spot spikes, one-
month futures spike almost 1:1 while two-year futures barely move.
A random-walk (Black--Scholes) spot model cannot reproduce this,
which is why Schwartz is preferred for commodities.
\end{intuitionbox}

\subsection{Half-life, stationarity, and what the parameters mean}

The log-spot's stationary distribution is
$X_\infty \sim \Normal\!\bigl(\alpha, \volat^2/(2\kappa)\bigr)$.
Two interpretable summary statistics follow.

\paragraph{Half-life of shocks.} If $X_0 = \alpha + \delta$, then
$\E[X_t - \alpha] = \delta e^{-\kappa t}$, half-life
$\tau_{1/2} = \ln 2 / \kappa$.  For $\kappa = 1$/yr,
$\tau_{1/2} \approx 8.3$ months; for electricity with
$\kappa \sim 50$/yr, $\tau_{1/2} \approx 5$ days.

\paragraph{Steady-state std of log-spot.}
$\sqrt{\volat^2/(2\kappa)}$ gives the one-sigma range around
$\alpha$.  For $\volat = 0.3$, $\kappa = 1$: $\pm 0.21$ ($\pm 21\%$
around equilibrium).  Large $\kappa$ means strong reversion and a
tight range.

\subsection{Risk-neutral drift and the market price of risk}

Switching $\Prob \to \Qm$ adjusts the long-run mean by the market
price of commodity risk $\lambda^{\mathrm{MPR}}$:
\[
  \alpha^* = \alpha - \frac{\lambda^{\mathrm{MPR}} \volat}{\kappa}.
\]
Risk-averse investors demand a premium for long commodity exposure,
so the $\Qm$ reversion target is \emph{lower} than the real-world
target.  This is the same Girsanov adjustment as in Black--Scholes,
applied to the mean-reversion target.  Calibration fits
$(\kappa, \alpha^*, \volat)$ by least-squares to the closed-form
curve (\ref{eq:schwartz_F}) against observed quotes.

\subsection{Extensions: two-factor and three-factor models}

The one-factor model has a critical flaw: $\alpha^*$ is constant,
so the long end of the curve never moves --- empirically it does,
slowly.  \citet{gibsonschwartz1990} make the convenience yield
$y_t$ a second mean-reverting factor, letting the curve twist.
\citet{schwartzsmith2000} recast this as a short-term mean-
reverting deviation plus a long-term equilibrium random walk (``S-S
two-factor''), often easier to calibrate.  Three-factor extensions
add stochastic rates.  For this book the one-factor model
suffices.
\citeSchwartz


%% ================================================================
\section{Commodity options on futures: Black's model}
\label{sec:black76_commodities}
%% ================================================================

Listed commodity options are almost universally \emph{options on
futures}.  A CL American call struck at \$80 delivers a CL futures
position on exercise, not a barrel.

\citet{black1976pricing} showed that options on futures can be
priced with the Black--Scholes machinery by substituting $F$ for
$\Spx$ and treating the futures as if it paid dividend yield
$\rrf$.  Intuition: a futures requires no up-front cash, so to make
the BS PDE balance one endows $F$ with dividend yield $\rrf$.

\begin{keyfactbox}[Black's formula for European options on futures]
For a European call and put on a futures with today's futures price
$F_0$, strike $\Kstrike$, maturity $\Tmat$, flat risk-free rate
$\rrf$, and futures volatility $\volat$,
\begin{align*}
  \Call &= e^{-\rrf \Tmat}\bigl[F_0 \CDFN(d_1) - \Kstrike \CDFN(d_2)\bigr], \\
  \Put  &= e^{-\rrf \Tmat}\bigl[\Kstrike \CDFN(-d_2) - F_0 \CDFN(-d_1)\bigr], \\
  d_{1,2} &= \frac{\ln(F_0/\Kstrike) \pm \tfrac12 \volat^2 \Tmat}{\volat \sqrt{\Tmat}}.
\end{align*}
The volatility $\volat$ is the lognormal vol of the futures price
itself, not of spot.
\end{keyfactbox}

\subsection{Worked example: 6-month WTI option}

Let $F_0 = \$82$, $\Kstrike = \$80$, $\Tmat = 0.5$ years,
$\volat = 30\%$, $\rrf = 5\%$.  Then
\begin{align*}
  d_1 &= \frac{\ln(82/80) + \tfrac12 (0.30)^2 (0.5)}{0.30 \sqrt{0.5}}
       = \frac{0.02469 + 0.0225}{0.21213} = 0.2225, \\
  d_2 &= d_1 - 0.30 \sqrt{0.5} = 0.0103.
\end{align*}
Reading $\CDFN(d_1) = 0.5880$, $\CDFN(d_2) = 0.5041$,
\begin{align*}
  \Call &= e^{-0.025}\bigl[82 \cdot 0.5880 - 80 \cdot 0.5041\bigr] \\
        &= 0.9753 \cdot \bigl[48.219 - 40.325\bigr] = 0.9753 \cdot 7.894
         = \$7.69.
\end{align*}
The put from put-call parity (which for options on futures reads
$\Call - \Put = e^{-\rrf T}(F_0 - K)$) gives
$\Put = 7.69 - e^{-0.025}(82-80) = 7.69 - 1.95 = \$5.74$.
(Verified numerically: $\Call = 7.6934$, $\Put = 5.7428$.)

\subsection{Python reference implementation}

\begin{lstlisting}[caption={Black's 1976 formula for European options on futures.}]
import math

def _phi(x):
    # standard normal CDF via erf
    return 0.5 * (1.0 + math.erf(x / math.sqrt(2.0)))

def black76(F0, K, T, sigma, r, kind="call"):
    """
    Price a European option on a futures contract (Black 1976).
      F0    : today's futures price
      K     : strike
      T     : time to expiry (years)
      sigma : futures lognormal vol (per-year)
      r     : continuously-compounded risk-free rate
      kind  : "call" or "put"
    """
    d1 = (math.log(F0 / K) + 0.5 * sigma * sigma * T) / (sigma * math.sqrt(T))
    d2 = d1 - sigma * math.sqrt(T)
    disc = math.exp(-r * T)
    if kind == "call":
        return disc * (F0 * _phi(d1) - K * _phi(d2))
    else:
        return disc * (K * _phi(-d2) - F0 * _phi(-d1))

# Worked example: 6M WTI call, F=82, K=80, sigma=30%, r=5%
print(black76(82, 80, 0.5, 0.30, 0.05, "call"))  # 7.6934
print(black76(82, 80, 0.5, 0.30, 0.05, "put"))   # 5.7428
\end{lstlisting}

\subsection{Options on spot vs options on futures: Samuelson's effect}

Not all commodity options are on futures.  LME metals options
reference the 3-month forward; OTC options on physical crude
reference a cash spot index.  For spot-referenced options the BS
formula uses $\rrf - y$ in place of the dividend yield --- the
Garman--Kohlhagen formula from \cref{ch:fx_derivatives} with $y$ as
the foreign rate.

For options on futures there is the \emph{Samuelson effect}:
futures vol rises as expiry approaches.  Long-dated futures are
anchored at the long-run mean; near-expiry futures converge to
spot and inherit its vol.  Under Schwartz,
\[
  \volat_F(t,T) = \volat\, e^{-\kappa(T-t)},
\]
growing from $\volat e^{-\kappa \Tmat}$ at option inception to
$\volat$ at futures expiry.  Feeding front-month vol into options
on long-dated futures systematically overprices them.

\begin{pitfallbox}[Don't plug spot vol into Black's model]
A frequent error calibrates $\volat$ from spot returns and feeds it
into Black's formula on futures.  These are different vols.  Under
Schwartz dynamics long-dated futures are \emph{less} volatile than
spot (mean-reversion anchors the long end to $\alpha^*$).  Always
calibrate $\volat$ from the futures return series you are pricing
against, or from implied vols of liquid options on that same
contract.
\end{pitfallbox}


%% ================================================================
\section{Seasonality}
\label{sec:seasonality}
%% ================================================================

Many commodities exhibit a periodic component in both their futures
curve and implied vol, repeating yearly.  This is a structural
feature of physical supply/demand cycles, not a statistical
artefact.

\begin{itemize}
  \item \textbf{Natural gas}: Northern winter heating demand lifts
    Dec--Feb Henry Hub futures above summer months; January is the
    winter peak, August the trough.  Storage injects all summer and
    withdraws all winter.
  \item \textbf{Agriculturals} (corn, soybeans, wheat): harvest
    brings a glut; old-crop (pre-harvest) trades at a premium to
    new-crop.  The \emph{crop year} creates a discontinuity at the
    harvest month.
  \item \textbf{Gasoline} (RBOB): summer driving demand and the
    costlier summer-grade blend peak in May--June.
  \item \textbf{Electricity}: diurnal and weekly seasonality on top
    of annual; peak vs off-peak blocks quoted separately.
\end{itemize}

To incorporate seasonality, write
\[
  \ln F(t,T) = \ln F_{\text{nonseasonal}}(t,T) + s(T),
\]
where $s(T)$ is a deterministic periodic function (a sum of
sinusoids fit to historical log-futures residuals after removing
curve-level effects).  Implied vol surfaces inherit the periodic
component --- winter gas options carry higher vol than summer.

\subsection{A Fourier-style seasonal overlay}

A simple parametric form uses annual and semi-annual sinusoids:
\[
  s(T) = a_1 \cos(2\pi T) + b_1 \sin(2\pi T)
       + a_2 \cos(4\pi T) + b_2 \sin(4\pi T),
\]
with $T$ in years.  The four coefficients are fitted by OLS to
log-futures residuals after stripping the Schwartz curve.  For
natural gas this captures most of the month-of-year effect; add a
third harmonic if residuals show a third peak.

\begin{keyfactbox}[Seasonality as a model layer]
Seasonality is treated as an \emph{additive deterministic layer} on
top of a non-seasonal stochastic model (Schwartz, Schwartz--Smith,
etc.). It does not require new random factors; it shifts the
curve's mean and the vol surface's level but preserves the
stochastic dynamics. Calibration: fit the periodic function to
residuals, then fit the stochastic model to the deseasonalised
series.
\end{keyfactbox}

\begin{pitfallbox}[Seasonality is not alpha]
A rookie mistake: January gas trades above August, so buy August /
sell January and declare victory.  The spread is priced into the
curve at all times; you are not earning a risk premium but paying
inter-seasonal storage on the short leg.  Seasonality is a
deterministic feature the market has already priced.  Tradeable
edges exist only in \emph{deviations} from the expected pattern
(unusually warm winter, wet harvest).
\end{pitfallbox}


%% ================================================================
\section{Physical vs financial settlement}
\label{sec:physical_financial_settlement}
%% ================================================================

A futures contract settles one of two ways.  \emph{Physical
settlement} delivers the underlying: the short delivers barrels,
the long pays the final settlement price times contract size.
\emph{Financial (cash) settlement} pays P\&L in cash against a
published reference and terminates without physical transfer.

\begin{itemize}
  \item \textbf{CL (WTI crude)}: physical, delivery at Cushing, OK.
    Long holders not closing out before expiry must take delivery of
    1{,}000 barrels.
  \item \textbf{BZ (Brent crude)}: financial, against ICE Brent
    Index (Platts).
  \item \textbf{NG (Henry Hub gas)}: physical, at the Henry Hub
    pipeline junction, Erath, LA.
  \item \textbf{HG (copper)}: physical, at LME warehouses.
  \item \textbf{Agriculturals (ZC, ZS, ZW)}: physical, at approved
    silos.
  \item \textbf{Power/electricity}: financial; no economic way to
    deliver 10 MWh into the grid on a specific day.
\end{itemize}

The distinction matters \emph{near expiry}.  A fund long physical-
settled oil must roll or arrange delivery (store, resell, or
transport).  The 2020 WTI event (20 April 2020, May contract at
$-\$37.63$) happened because long holders unable to take delivery
paid shorts to take the contract away.

\begin{intuitionbox}
Financial settlement is a bet on the index; physical is a promise
to move real stuff.  Physical settlement creates a hard deadline
for speculators: be flat or hire a tanker.  Fund-level commodity
exposure is held through index replication or rolled front-month
futures, with rolls executed mechanically 5--10 business days
before First Notice Day to avoid physical assignment.
\end{intuitionbox}


%% ================================================================
\section{Case studies from history}
\label{sec:commodity_history}
%% ================================================================

\begin{historybox}[title={Schwartz 1997: the canonical reference}]
\citet{schwartz1997stochastic} compared three stochastic models for
commodity spot: a one-factor mean-reverting model (``Schwartz 1''),
a two-factor with stochastic convenience yield, and a three-factor
adding stochastic rates.  The paper is why commodity options desks
still structure pricing libraries around mean-reverting log-spot
plus optional stochastic $y_t$.  Schwartz later collaborated with
Smith on the short-term/long-term reformulation that became the
industry standard.
\end{historybox}

\begin{historybox}[title={Metallgesellschaft 1993: stack-and-roll blowup}]
Metallgesellschaft AG sold long-dated fixed-price oil contracts to
US customers and hedged by \emph{stack-and-roll}: a large stack of
near-month futures rolled forward each month.  When WTI flipped
from backwardation to steep contango in late 1993, the roll turned
punitive (sell cheap expiring, buy expensive next month).
Accounting forced MG to mark the futures leg while leaving the
offsetting long-dated customer contracts unmarked; the board
liquidated the hedge at the worst moment, crystallising $\sim$\$1.3
billion of loss.  Lesson: a theoretically correct hedge can destroy
a firm when funding cannot survive the interim cash outflows.
\end{historybox}

\begin{historybox}[title={Amaranth 2006: \$6B on a gas-spread bet}]
Amaranth Advisors, via trader Brian Hunter, built an enormous long
position in winter gas calendar spreads (long March, short April),
betting winter scarcity would widen them.  When spreads collapsed
in September 2006 the fund lost $\sim$\$6.5 billion in a week and
closed.  Their position was $\sim 40\%$ of NYMEX March--April open
interest --- they were both the market and its victim.  Lesson:
seasonality trades are seductive but already priced, edges require
genuine divergence forecasts, and size must stay a fraction of
open interest.
\end{historybox}


%% ================================================================
\section{Commodity playbooks}
\label{sec:commodity_playbooks}
%% ================================================================

\begin{prosperitybox}[title={Schwartz-like dynamics in Prosperity}]
KELP (Prosperity 3) and its successors show the Schwartz one-factor
signature: mean-reverting log-spot with a few-hundred-tick half-
life and a well-defined long-run mean.  The Frankfurt Hedgehogs'
winning KELP approach estimated $\kappa$ and $\alpha$ from a
rolling sample and skewed quotes toward $\hat\alpha$: buy below,
sell above.  This is a textbook Schwartz-style trade: mean-revert
log-spot, collect spread on noise-trader deviations, flatten
inventory as price reverts.  Recognise a mean-reverting log-spot
and the Schwartz toolkit applies --- no commodities needed.
\end{prosperitybox}

\begin{gsbox}[title={Commodities Strats on the trading floor}]
A Goldman commodities Strat spends the day on four things.
(i) \emph{Curve construction}: bootstrapping daily futures curves
for WTI, Brent, Henry Hub, corn, soy, and power hubs; QC against
stale quotes, rolls, and holiday discontinuities.
(ii) \emph{Schwartz or Schwartz--Smith calibration}: refitting
$(\kappa, \alpha^*, \volat)$ each morning, with overrides on
inventory data, OPEC, or weather.
(iii) \emph{Options-on-futures pricer}: a Black--76 library with
American exercise, vol surface interpolation (seasonally adjusted),
and Greeks including the commodity-specific
$\partial V / \partial y$.
(iv) \emph{Seasonality} overlay, recalibrated quarterly.  All four
auto-run overnight; the trader inspects a curve-diff summary at the
7am run-through.
\end{gsbox}


%% ================================================================
\section{Key facts to memorise}
\label{sec:commodity_keyfacts}
%% ================================================================

\begin{itemize}
  \item \textbf{Commodity cost-of-carry:} for a storable commodity,
    $F(T) = \Spx e^{(\rrf + u - y)T}$, where $u$ is the storage cost
    rate and $y$ is the convenience yield.
  \item \textbf{Contango} = upward-sloping curve ($F(T_2) > F(T_1)$
    for $T_2 > T_1$); typical of loose inventories.
    \textbf{Backwardation} = downward-sloping; typical of scarcity.
  \item \textbf{Contango is not a price forecast.} $F(T)$ is a
    risk-adjusted forward price, not $\E[\Spx_T]$.  A contangoed
    market can still see spot fall.
  \item \textbf{Implied convenience yield}
    $y(T) = \rrf(T) + u(T) - T^{-1}\ln(F(T)/\Spx)$ backs out the
    per-barrel scarcity signal from observable quotes.
  \item \textbf{Schwartz one-factor:} log-spot is OU with reversion
    speed $\kappa$ and long-run log-mean $\alpha$.  The long end of
    the futures curve is anchored at a constant equilibrium level
    $\exp(\alpha^* + \volat^2/(4\kappa))$.
  \item \textbf{Schwartz futures price:} given in closed form by
    equation (\ref{eq:schwartz_F}).  Long-maturity futures do not
    move 1:1 with spot.
  \item \textbf{Black 1976:} price options on futures by the
    Black--Scholes formula with $F$ in place of $\Spx$ and an
    extra discount factor $e^{-\rrf T}$ in front.  Use futures
    vol, not spot vol.
  \item \textbf{Seasonality} is a deterministic periodic overlay on
    top of the stochastic model, both for the curve and for the
    vol surface.  Tradeable edges are \emph{deviations} from the
    pattern, not the pattern itself.
  \item \textbf{Physical vs financial settlement:} CL/NG/agris
    settle physically; Brent, power, and most vol products settle
    financially.  Physical settlement matters near expiry --- roll
    or be ready for delivery.
  \item \textbf{Metallgesellschaft, Amaranth:} a mathematically
    correct hedge can destroy a firm if the curve's shape moves
    against the funding path.  Always stress-test the
    cash-outflow profile, not just the terminal P\&L.
\end{itemize}
